\chapter{Numbers, Alignment, and Type-Trait Additions}
\label{ch:c20-numbers}

\begin{quote}
\emph{The final cluster of C++20 library additions serves numeric and generic code: \texttt{<numbers>} provides typed mathematical constants; \texttt{std::midpoint} and \texttt{std::lerp} give overflow-safe averaging and interpolation; \texttt{std::assume\_aligned} promises the optimizer a pointer's alignment; and new type traits (\texttt{remove\_cvref}, \texttt{type\_identity}, \texttt{is\_bounded\_array}, \texttt{common\_reference}) refine metaprogramming.}
\end{quote}

\section{The <numbers> Header: Mathematical Constants}
\label{sec:c20-numbers-header}

\begin{lstlisting}[language=C++,caption={Typed, full-precision mathematical constants}]
#include <numbers>

double pi   = std::numbers::pi;          // 3.14159265358979... (double)
float  pi_f = std::numbers::pi_v<float>; // float-precision pi
double e    = std::numbers::e;
double sqrt2= std::numbers::sqrt2;
double phi  = std::numbers::phi;         // golden ratio
double area = std::numbers::pi * r * r;  // no #define, no magic digits
\end{lstlisting}

\begin{center}
\begin{tabular}{ll}
\hline
\textbf{Constant} & \textbf{Value} \\
\hline
\texttt{pi} & $\pi$ \\
\texttt{e} & Euler's number \\
\texttt{sqrt2}, \texttt{sqrt3} & $\sqrt2$, $\sqrt3$ \\
\texttt{ln2}, \texttt{ln10} & natural logs \\
\texttt{phi} & golden ratio \\
\texttt{inv\_pi}, \texttt{inv\_sqrt2} & reciprocals \\
\texttt{egamma} & Euler--Mascheroni constant \\
\hline
\end{tabular}
\end{center}

The \texttt{\_v<T>} forms give the constant at a chosen precision; all are \texttt{constexpr}.

\section{std::midpoint: Overflow-Safe Averaging}
\label{sec:c20-midpoint}

\begin{lstlisting}[language=C++,caption={Averaging without overflow}]
#include <numeric>
#include <cstdint>

std::int32_t a = 2'000'000'000, b = 2'000'000'000;
// int mid_bad = (a + b) / 2;            // overflow -> UB
std::int32_t mid = std::midpoint(a, b);  // 2'000'000'000, computes a + (b-a)/2

int arr[100];
int* m = std::midpoint(arr, arr + 100);  // arr + 50, no pointer overflow
\end{lstlisting}

Fixes the famous binary-search overflow bug; overloaded for integers, floats, and pointers.

\section{std::lerp: Linear Interpolation}
\label{sec:c20-lerp}

\begin{lstlisting}[language=C++,caption={Well-behaved linear interpolation}]
#include <cmath>

double a = 0.0, b = 10.0;
double mid   = std::lerp(a, b, 0.5);   // 5.0
double start = std::lerp(a, b, 0.0);   // exactly a
double end   = std::lerp(a, b, 1.0);   // exactly b (guaranteed)
double beyond= std::lerp(a, b, 2.0);   // 20.0 — extrapolation
\end{lstlisting}

Guarantees exact endpoints and monotonicity in \texttt{t}, unlike the naive \texttt{a + t*(b-a)}; \texttt{constexpr}.

\section{std::assume\_aligned: Promising Alignment to the Optimizer}
\label{sec:c20-assume-aligned}

\begin{lstlisting}[language=C++,caption={Informing the optimizer of alignment}]
#include <memory>
#include <cstddef>

void scale(float* data, std::size_t n, float k) {
    float* aligned = std::assume_aligned<64>(data);   // promise 64-byte alignment
    for (std::size_t i = 0; i < n; ++i)
        aligned[i] *= k;          // may emit aligned SIMD loads/stores
}
\end{lstlisting}

A hard promise: if the pointer is not actually aligned to N, behavior is undefined. Use only when you control the allocation.

\section{Type-Trait Additions: remove\_cvref and type\_identity}
\label{sec:c20-remove-cvref}

\begin{lstlisting}[language=C++,caption={remove\_cvref and type\_identity}]
#include <type_traits>

static_assert(std::is_same_v<std::remove_cvref_t<const int&>, int>);

template <typename T>
void store(T&& value) {
    using Stored = std::remove_cvref_t<T>;   // bare value type for a forwarding ref
    Stored copy = std::forward<T>(value);
}

template <typename T>
void clamp_to(T value, std::type_identity_t<T> lo, std::type_identity_t<T> hi);
// T deduced ONLY from 'value'; bounds are a non-deduced context.
\end{lstlisting}

\section{Type-Trait Additions: is\_bounded\_array and common\_reference}
\label{sec:c20-is-bounded-array}

\begin{lstlisting}[language=C++,caption={Array classification and common\_reference}]
#include <type_traits>

static_assert(std::is_bounded_array_v<int[10]>);   // true  — known size
static_assert(std::is_unbounded_array_v<int[]>);   // true  — unknown size

using CR = std::common_reference_t<int&, const int&>;   // const int&
static_assert(std::is_same_v<CR, const int&>);
\end{lstlisting}

\texttt{common\_reference} underpins the Ranges/iterator concepts (relating an iterator's \texttt{reference} and \texttt{value\_type}).

\section{Professional Insights}
\label{sec:c20-numbers-insights}

\textbf{Replace \texttt{M\_PI} and hand-typed literals with \texttt{<numbers>}} --- standard, full-precision, \texttt{constexpr}, precision-parameterized; eliminates the seed-precision bug.

\textbf{Use \texttt{std::midpoint} for every bisection} --- fixes the famous \texttt{(low+high)/2} overflow; overloaded for pointers.

\textbf{Reach for \texttt{std::lerp} whenever interpolation feeds further computation} --- exact endpoints and monotonicity the naive formula lacks.

\textbf{Treat \texttt{std::assume\_aligned} as a hard, dangerous promise} --- UB if the alignment is wrong; use only on allocations you control.

\textbf{Adopt \texttt{std::remove\_cvref\_t} as the standard forwarding-reference value trait} --- and know the niches of \texttt{type\_identity}, \texttt{is\_bounded\_array}, and \texttt{common\_reference}.
